\section{实验结果与分析}

\subsection{数据集与评价指标}

实验选用KITTI跟踪数据集中的11个序列（涵盖城市道路、居民区与高速路段）以及自采集的北方某城市10个交叉口的实际轨迹数据。所有轨迹经人工标注后作为真值。评估指标包括：均方根误差（RMSE）、平均绝对误差（MAE）、轨迹平滑性（Ts）以及滤波一致性（NCI）。其中轨迹平滑性定义为相邻位移向量的夹角变化方差，滤影一致性反映新息序列的白化程度。

\subsection{定量对比}

表\ref{tab:quantitative}列出不同滤波器在测试集上的性能平均值。本方法在全部四项指标上均取得最优：RMSE较标准卡尔曼滤波（KF）下降15.3%，较扩展卡尔曼滤波（EKF）下降10.1%；MAE分别下降14.2%与9.7%；轨迹平滑性提升幅度分别为21.6%与12.2%；滤波一致性指标也反映了模型估计与实际不确定度的高度匹配。

\begin{table}[htbp]
    \centering
    \caption{各方法在测试集上的定量性能对比}
    \label{tab:quantitative}
    \begin{tabular}{lcccc}
        \hline
        方法 & RMSE (m) & MAE (m) & 平滑性 (rad²) & NCI \\
        \hline
        KF & 0.782 & 0.541 & 0.0338 & 0.91 \\
        EKF & 0.737 & 0.514 & 0.0302 & 0.93 \\
        Ours & 0.663 & 0.464 & 0.0265 & 0.97 \\
        \hline
    \end{tabular}
\end{table}

\subsection{定性分析}

图\ref{fig:trajectory}展示了一个典型遮挡场景下的轨迹复原效果。在观测出现约5.0秒的连续缺失区间（如图中灰色区域）内，本方法能够保持与真值高度吻合的预测，而标准KF在该时段出现明显漂移。图\ref{fig:error_dist}进一步给出了误差分布的箱线图，显示本方法的误差上分位数显著小于对比方法，表明其在极端工况下的鲁棒性更强。